% ============================================================
% f2013_shanghai_science_q13_rotation.tex
% 2013 上海卷（理科）第 13 题 区域 D 与绕 y 轴旋转示意（TikZ 矢量重绘）
% 边界：圆 (x-1)^2+y^2=1 的右半支、圆 (x-3)^2+y^2=1 的右半支，
%       以及直线 y=1、y=-1。阴影部分为 D。
% ============================================================
\begin{tikzpicture}[
  x=1.15cm, y=1.35cm,
  every node/.style={font=\small, inner sep=1pt},
  line/.style={line width=0.6pt},
  axis/.style={line width=0.5pt, -{Stealth[length=2.2mm,width=1.6mm]}},
]

  % 阴影区域 D
  \fill[black!12]
    (1,1) -- (3,1)
    arc[start angle=90, end angle=-90, radius=1]
    -- (1,-1)
    arc[start angle=-90, end angle=90, radius=1]
    -- cycle;

  % D 的边界
  \draw[line]
    (1,1) -- (3,1)
    arc[start angle=90, end angle=-90, radius=1]
    -- (1,-1)
    arc[start angle=-90, end angle=90, radius=1]
    -- cycle;

  % 坐标轴
  \draw[axis] (-0.35,0) -- (4.55,0) node[below] {$x$};
  \draw[axis] (0,-1.45) -- (0,1.55) node[left] {$y$};

  % x 轴刻度
  \foreach \x in {1,2,3,4}{
    \draw[line] (\x,0.05) -- (\x,-0.05);
    \node[below] at (\x,-0.07) {$\x$};
  }
  % y 轴刻度
  \foreach \y in {-1,1}{
    \draw[line] (-0.05,\y) -- (0.05,\y);
    \node[left] at (-0.07,\y) {$\y$};
  }

  \node at (-0.17,-0.18) {$O$};

\end{tikzpicture}
